% !PRINTED 587
\weberpage{587}
On the other hand, $f_\nu(\tau)\,\psi_\mu(\sigma)$ is an integral functional of the field
$\Omega$, and consequently, by the theorem \S.\,158, 4. already applied above:

\weq{16}{
\begin{aligned}
S_\Omega \left( \frac{f_\nu(\tau)\,\psi_\mu(\sigma)}{F'(\tau)} \right) &= b_{\mu,\nu} &
\mu &= 0, 1, \ldots, m-1\\
& & \nu &= 0, 1, \ldots, n-1
\end{aligned}
}
is an integral rational functional. If we multiply this equation by
$\tau^{n-\nu-1}\sigma^{m-\mu-1}$ and take the sum over all $\mu$ and $\nu$, then,
as before — since for two roots $\sigma_k, \tau_i$ of $\psi=0, f=0$ different
from $\sigma, \tau$
\[
\begin{aligned}
\sigma^{m-1}\psi_0(\sigma_k) + \sigma^{m-2}\psi_1(\sigma_k) + \cdots + \psi_{m-1}(\sigma_k) &= 0,\\
\tau^{n-1} f_0(\tau_i) + \tau^{n-2} f_1(\tau_i) + \cdots + f_{n-1}(\tau_i) &= 0,
\end{aligned}
\]
and
\[
\begin{aligned}
\sigma^{m-1}\psi_0(\sigma) + \sigma^{m-2}\psi_1(\sigma) + \cdots + \psi_{m-1}(\sigma) &= \psi'(\sigma),\\
\tau^{n-1} f_0(\tau) + \tau^{n-2} f_1(\tau) + \cdots + f_{n-1}(\tau) &= f'(\tau)
\end{aligned}
\]
hold — we obtain the equation
\weq{17}{
\frac{f'(\tau)\,\psi'(\sigma)}{F'(\tau)} = \sum^{\mu,\nu} b_{\mu,\nu}\, \tau^{n-\nu-1}\sigma^{m-\mu-1}.
}

Hence $f'(\tau)\,\psi'(\sigma)$ too is divisible by $F'(\tau)$, and our Theorem (8)
is therefore proved.

If we denote by $G_\Omega, G_R$ the fundamental ideals of the fields $\Omega, R$, and by
$\mathfrak{G}_R$ the relative fundamental ideal of $\Omega$ with respect to $R$, then the
theorem just proved may also be given the following expression:
\weq{18}{
G_\Omega = \mathfrak{G}_R\, G_R.
}
It is clear how this theorem may be generalised: Let
\[
\Omega_1, \Omega_2, \Omega_3, \ldots
\]
be a sequence of fields, each of which contains the following ones as divisors. Further, let
$\mathfrak{G}_{i,k}$ be the fundamental ideal of $\Omega_i$ with respect to
$\Omega_k$ ($k > i$); then
\weq{19}{
\mathfrak{G}_{i,k} = \mathfrak{G}_{i,i+1}\,\mathfrak{G}_{i+1,i+2} \cdots \mathfrak{G}_{k-1,k}.
}
The deeper significance of these fundamental ideals will emerge still more clearly in one
of the following sections.

But we wish to call attention already here to a remarkable phenomenon: namely, that it can
happen that $f'(\tau)$ is a unit, which by \gesperrt{Minkowski's} theorem is not
possible when $R$ is the field of rational numbers. In these cases, of which the theory of com-
%CONT>>
